
\documentclass[11pt,a4paper]{article}

\usepackage[a4paper,margin=2.35cm,headheight=24pt]{geometry}
\usepackage[french]{babel}
\usepackage{lmodern}
\usepackage[table]{xcolor}
\usepackage{graphicx}
\usepackage{booktabs,longtable,tabularx,array}
\usepackage{caption,subcaption,float}
\usepackage{hyperref}
\usepackage{bookmark}
\usepackage{fancyhdr}
\usepackage{lastpage}
\usepackage{enumitem}
\usepackage[most]{tcolorbox}
\usepackage{fvextra}
\usepackage{microtype}
\usepackage{parskip}
\usepackage{titlesec}
\usepackage{ragged2e}
\usepackage{needspace}
\usepackage{tocloft}
\usepackage{etoolbox}

\definecolor{SectionBlue}{HTML}{1F4E8C}
\definecolor{BoxBlueBg}{HTML}{F5F7FA}
\definecolor{BoxOrangeBg}{HTML}{FFF5EB}
\definecolor{BoxOrange}{HTML}{E67E22}
\definecolor{BoxGreenBg}{HTML}{EEF8F0}
\definecolor{BoxGreen}{HTML}{2E8B57}
\definecolor{SoftGray}{HTML}{F6F6F6}
\definecolor{GrayLight}{HTML}{F6F8FA}
\definecolor{RuleGray}{HTML}{CFCFCF}
\definecolor{CodeBg}{HTML}{F5F5F5}
\definecolor{CodeFrame}{HTML}{BDBDBD}

\hypersetup{
  colorlinks=true,
  linkcolor=SectionBlue,
  urlcolor=SectionBlue,
  pdftitle={Rapport detaille de realisation - Seance 3 MongoDB},
  pdfauthor={OpenAI},
  pdfsubject={Rapport technique avec preuves d'execution},
  pdfcreator={LuaLaTeX}
}

\pagestyle{fancy}
\fancyhf{}
\lhead{\textcolor{SectionBlue}{Rapport de réalisation --- MongoDB / Séance 3}}
\rhead{\textcolor{SectionBlue}{\thepage}}
\renewcommand{\headrulewidth}{0.4pt}
\renewcommand{\headrule}{\hbox to\headwidth{\color{black}\leaders\hrule height \headrulewidth\hfill}}
\cfoot{}

\fancypagestyle{plain}{
  \fancyhf{}
  \lhead{\textcolor{SectionBlue}{Rapport de réalisation --- MongoDB / Séance 3}}
  \rhead{\textcolor{SectionBlue}{\thepage}}
  \renewcommand{\headrulewidth}{0.4pt}
  \renewcommand{\headrule}{\hbox to\headwidth{\color{black}\leaders\hrule height \headrulewidth\hfill}}
  \cfoot{}
}

\renewcommand{\contentsname}{Table des matières}
\setlength{\parindent}{0pt}
\setlength{\parskip}{0.55em}
\setlength{\cftbeforesecskip}{7pt}
\setlength{\cftbeforesubsecskip}{2pt}
\renewcommand{\cftsecfont}{\color{SectionBlue}\bfseries}
\renewcommand{\cftsecpagefont}{\bfseries}
\renewcommand{\cftsubsecfont}{\color{SectionBlue}}
\renewcommand{\cfttoctitlefont}{\Huge\bfseries\color{SectionBlue}}
\renewcommand{\cftaftertoctitle}{\hfill}

\titleformat{\section}{\Large\bfseries\color{SectionBlue}}{\thesection}{0.75em}{}
\titleformat{\subsection}{\large\bfseries\color{SectionBlue}}{\thesubsection}{0.75em}{}
\titleformat{\subsubsection}{\normalsize\bfseries\color{SectionBlue}}{\thesubsubsection}{0.75em}{}
\titlespacing*{\section}{0pt}{1.3em}{0.7em}
\titlespacing*{\subsection}{0pt}{1.0em}{0.45em}
\titlespacing*{\subsubsection}{0pt}{0.9em}{0.35em}

\setlist[itemize]{topsep=2pt,itemsep=2pt,leftmargin=1.5em}
\setlist[enumerate]{topsep=2pt,itemsep=2pt,leftmargin=1.7em}
\newcolumntype{Y}{>{\RaggedRight\arraybackslash}X}
\newcolumntype{C}[1]{>{\Centering\arraybackslash}p{#1}}
\renewcommand{\arraystretch}{1.16}

\DeclareCaptionLabelSeparator{endash}{\space--\space}
\captionsetup{font=small,labelfont=sc,labelsep=endash}

\tcbset{enhanced,breakable,boxrule=0.75pt,arc=8pt,left=10pt,right=10pt,top=8pt,bottom=8pt}

\newtcolorbox{summarybox}[1][]{%
  colback=BoxBlueBg,
  colframe=SectionBlue,
  before upper={\ifstrempty{#1}{}{\textbf{#1}.\ }},
}

\newtcolorbox{warnbox}[1][]{%
  colback=BoxOrangeBg,
  colframe=BoxOrange,
  before upper={\ifstrempty{#1}{}{\textbf{#1}.\ }},
}

\newtcolorbox{resultbox}[1][]{%
  colback=BoxGreenBg,
  colframe=BoxGreen,
  before upper={\ifstrempty{#1}{}{\textbf{#1}.\ }},
}

\newcommand{\statusok}{\textcolor{BoxGreen}{\textbf{Oui}}}
\newcommand{\statuspartial}{\textcolor{BoxOrange}{\textbf{Partiel}}}
\newcommand{\statusno}{\textcolor{red!70!black}{\textbf{Non démontré}}}

\newcommand{\evidencepage}[3]{%
\begin{figure}[H]
\centering
\includegraphics[width=0.90\linewidth,height=0.68\textheight,keepaspectratio]{evidence/#1}
\caption{#2}
\end{figure}
\noindent\textbf{Commentaire d'analyse.} #3
\clearpage
}

\begin{document}
\pagenumbering{gobble}

\begin{titlepage}
\centering
\vspace*{1.6cm}

{\fontsize{28}{34}\selectfont\bfseries\color{SectionBlue}Rapport détaillé de réalisation\par}
\vspace{0.55cm}
{\fontsize{18}{22}\selectfont Module NoSQL Security --- Séance 3\par}
\vspace{0.35cm}
{\Large Authentification, RBAC, exposition réseau et hygiène de configuration\par}

\vspace{1.0cm}
\begin{summarybox}[Objet du document]
Ce rapport reconstruit le travail effectivement réalisé à partir du support de séance en HTML et de l'archive de captures d'écran fournie. L'objectif est de produire un livrable exploitable tel quel, rédigé en LaTeX / LuaLaTeX, propre à l'impression et suffisamment détaillé pour documenter l'exécution des travaux pratiques, les tests autorisés / refusés et l'état réel du hardening observé.
\end{summarybox}

\vfill
\begin{minipage}{0.78\textwidth}
\centering
\begin{tabular}{>{\bfseries}p{4.8cm}p{7.2cm}}
Support pédagogique & \texttt{seance3\_nosql\_security\_Partie3(2).html} \\
Preuves visuelles & \texttt{Lecture 3 Screenshots.zip} (16 captures) \\
Base de travail & \texttt{soc\_lite} \\
Environnement visé & Docker + MongoDB 7 / mongosh \\
Livrables produits & PDF + code source LuaLaTeX \\
\end{tabular}
\end{minipage}

\vfill
\end{titlepage}

\clearpage
\pagenumbering{arabic}
\setcounter{page}{1}
\tableofcontents
\clearpage
\section{Objet du rapport}

Le présent rapport documente le travail réalisé dans le cadre de la \textbf{Séance 3 -- Hardening MongoDB}. L'objectif de la séance était de démontrer, sur une base NoSQL de type SOC, la différence entre une instance MongoDB exposée sans authentification et une instance correctement durcie par :
\begin{itemize}
  \item l'activation de l'authentification,
  \item la création d'un administrateur initial,
  \item la gestion correcte de la base d'authentification (\texttt{authenticationDatabase}),
  \item la mise en place d'un contrôle d'accès par rôles (RBAC),
  \item la validation explicite des opérations \emph{autorisées} et \emph{refusées}.
\end{itemize}

Le support pédagogique attendait un mini-rapport ``avant / après'' ainsi qu'une checklist de hardening. Ce document remplit cette finalité et ajoute une analyse de cohérence entre la politique visée et les résultats réellement observés.

\section{Sources et méthode de reconstruction}

\subsection{Corpus exploité}

Deux sources ont été exploitées :
\begin{enumerate}
  \item le support HTML de la séance, qui définit les objectifs, les activités, les commandes attendues et les contrôles à vérifier ;
  \item les 16 captures d'écran fournies, qui permettent de reconstituer chronologiquement les manipulations effectivement exécutées.
\end{enumerate}

\subsection{Approche d'analyse}

La reconstruction a été menée selon la logique suivante :
\begin{itemize}
  \item ordonner les captures par horodatage ;
  \item identifier, pour chaque capture, la commande exécutée, le contexte MongoDB visé et le résultat obtenu ;
  \item regrouper les preuves en trois blocs : \textbf{exposition insecure}, \textbf{authentification}, \textbf{RBAC / moindre privilège} ;
  \item comparer les résultats observés avec la cible de sécurité du support ;
  \item signaler les écarts significatifs, notamment la limite des rôles built-in \texttt{read} et \texttt{readWrite} sur la collection \texttt{users}.
\end{itemize}
\newpage
\begin{summarybox}[Synthèse immédiate]
Le jeu de preuves fourni documente de manière convaincante l'exécution des trois activités principales du support :
\begin{enumerate}
  \item démonstration d'un MongoDB \emph{insecure} ;
  \item activation et validation de l'authentification ;
  \item conception et test d'une politique RBAC.
\end{enumerate}
En revanche, la \textbf{correction finale du devoir maison} -- à savoir la création de rôles personnalisés supplémentaires pour empêcher l'accès à \texttt{users} aux comptes \texttt{soc\_app} et \texttt{soc\_analyst\_ro} -- n'apparaît pas comme exécutée dans les captures. Cette partie reste donc \textbf{à compléter}.
\end{summarybox}

\section{Environnement et périmètre technique}

Le travail est réalisé dans un environnement Docker avec \texttt{mongosh}. La base cible est \texttt{soc\_lite}. Les objets manipulés au cours des démonstrations sont :
\begin{itemize}
  \item \texttt{security\_logs} : journaux de sécurité,
  \item \texttt{incidents} : incidents SOC,
  \item \texttt{users} : référentiel utilisateurs.
\end{itemize}

Trois profils fonctionnels sont ensuite modélisés :
\begin{itemize}
  \item \texttt{soc\_app} : compte applicatif de collecte,
  \item \texttt{soc\_analyst\_ro} : analyste en lecture,
  \item \texttt{soc\_manager} : responsable incidents.
\end{itemize}

\section{Chronologie détaillée des travaux réalisés}

\subsection{Phase 1 -- Démonstration d'une instance MongoDB insecure}

La première séquence de captures montre le lancement d'une instance MongoDB sans authentification et avec exposition réseau large. Cette phase répond à l'objectif pédagogique de visualiser le risque d'un service accessible sans contrôle.

\begin{tcolorbox}[breakable,colback=CodeBg,colframe=CodeFrame,title=Extrait de commandes observées]
\begin{Verbatim}[fontsize=\small,breaklines,breakanywhere]
docker run -d --name mongo_insecure -p 27017:27017 mongo:7
docker exec -it mongo_insecure mongosh

use soc_lite
db.security_logs.find().limit(3)
db.incidents.drop()
db.dropDatabase()
\end{Verbatim}
\end{tcolorbox}

\subsubsection*{Constats techniques}
\begin{itemize}
  \item Le port \texttt{27017} est publié sans restriction d'interface via \texttt{-p 27017:27017}.
  \item La connexion à \texttt{mongosh} est réalisée \textbf{sans identifiants}.
  \item Des opérations destructrices sont acceptées : suppression de la collection \texttt{incidents}, puis suppression de la base \texttt{soc\_lite}.
\end{itemize}

\begin{resultbox}[Résultat de sécurité]
Cette séquence établit clairement qu'une instance MongoDB démarrée sans authentification autorise la lecture, l'écriture et même la destruction de données par tout client ayant accès au port d'écoute.
\end{resultbox}

\subsubsection*{Preuve complémentaire d'écriture et d'administration anonymes}

Une seconde capture confirme que l'accès anonyme permet non seulement la lecture mais aussi l'insertion et l'administration du serveur :

\begin{tcolorbox}[breakable,colback=CodeBg,colframe=CodeFrame,title=Commandes observées]
\begin{Verbatim}[fontsize=\small,breaklines,breakanywhere]
use soc_lite
db.security_logs.insertOne({
  timestamp: new Date(),
  source_ip: "10.0.0.99",
  event_type: "test_insecure",
  message: "Ce document a été inséré sans authentification"
})

db.security_logs.find({ event_type: "test_insecure" })
db.adminCommand({ listDatabases: 1 })
\end{Verbatim}
\end{tcolorbox}

Le retour \texttt{acknowledged: true} ainsi que la réussite de \texttt{listDatabases} prouvent qu'un client non authentifié possède un contrôle complet sur l'instance.

\subsection{Phase 2 -- Activation de l'authentification et validation du compte administrateur}

La séquence suivante montre le basculement vers une instance sécurisée.

\begin{tcolorbox}[breakable,colback=CodeBg,colframe=CodeFrame,title=Commandes observées (version sécurisée)]
\begin{Verbatim}[fontsize=\small,breaklines,breakanywhere]
docker rm -f mongo_insecure

docker run -d --name mongo_secure \
  -p 127.0.0.1:27017:27017 \
  -e MONGO_INITDB_ROOT_USERNAME=<admin> \
  -e MONGO_INITDB_ROOT_PASSWORD=<mot_de_passe_admin> \
  mongo:7

docker exec -it mongo_secure mongosh
use soc_lite
db.security_logs.find()
\end{Verbatim}
\end{tcolorbox}

\subsubsection*{Constats techniques}
\begin{itemize}
  \item Le port est cette fois lié à \texttt{127.0.0.1}, ce qui réduit l'exposition réseau côté hôte.
  \item Toute opération sur les données retourne \texttt{MongoServerError[Unauthorized]: Command find requires authentication}.
  \item La protection n'est donc plus théorique : elle est validée par un refus explicite.
\end{itemize}

L'administrateur initial est ensuite utilisé avec succès :

\begin{tcolorbox}[breakable,colback=CodeBg,colframe=CodeFrame,title=Validation du compte administrateur]
\begin{Verbatim}[fontsize=\small,breaklines,breakanywhere]
docker exec -it mongo_secure mongosh -u <admin> -p <mot_de_passe_admin>
db.adminCommand({ listDatabases: 1 })
\end{Verbatim}
\end{tcolorbox}

Le résultat montre la liste des bases disponibles, ce qui confirme :
\begin{enumerate}
  \item que l'authentification est activée ;
  \item que le compte administrateur initial fonctionne ;
  \item que la protection distingue désormais clairement un accès anonyme d'un accès authentifié.
\end{enumerate}

\subsection{Phase 3 -- Validation de la création manuelle de l'administrateur initial}

Une autre séquence montre une variante intéressante : la création manuelle de l'administrateur sur un conteneur lancé avec \texttt{mongod --auth --bind\_ip\_all}. Cette partie démontre la bonne compréhension du mécanisme de \emph{localhost exception}.

\begin{tcolorbox}[breakable,colback=CodeBg,colframe=CodeFrame,title=Séquence manuelle observée]
\begin{Verbatim}[fontsize=\small,breaklines,breakanywhere]
docker run -d --name mongo_manual -p 127.0.0.1:27017:27017 mongo:7 \
  mongod --auth --bind_ip_all

docker exec -it mongo_manual mongosh
use admin
db.createUser({
  user: "<admin>",
  pwd: "<mot_de_passe_admin>",
  roles: [{ role: "root", db: "admin" }]
})

quit()
docker exec -it mongo_manual mongosh -u <admin> -p <mot_de_passe_admin>
db.adminCommand({ listDatabases: 1 })
\end{Verbatim}
\end{tcolorbox}

\begin{resultbox}[Intérêt de cette étape]
Cette démonstration va au-delà de la simple utilisation des variables d'environnement Docker : elle montre que l'étudiant a également validé le scénario pédagogique de création manuelle du premier compte root.
\end{resultbox}

\subsection{Phase 4 -- Validation de l'authenticationDatabase}

La notion de \texttt{authenticationDatabase} est explicitement testée dans les captures. Un compte \texttt{soc\_app} est utilisé avec la bonne base d'authentification, puis avec une base erronée.

\begin{tcolorbox}[breakable,colback=CodeBg,colframe=CodeFrame,title=Tests observés]
\begin{Verbatim}[fontsize=\small,breaklines,breakanywhere]
docker exec -it mongo_manual mongosh -u soc_app -p <mot_de_passe_app> \
  --authenticationDatabase soc_lite

use soc_lite
db.security_logs.find()
\end{Verbatim}
\end{tcolorbox}

La connexion réussit lorsque la base d'authentification est \texttt{soc\_lite}.

\begin{tcolorbox}[breakable,colback=CodeBg,colframe=CodeFrame,title=Erreur observée avec mauvaise authDB]
\begin{Verbatim}[fontsize=\small,breaklines,breakanywhere]
docker exec -it mongo_manual mongosh -u soc_app -p <mot_de_passe_app> \
  --authenticationDatabase admin

MongoServerError: Authentication failed.
\end{Verbatim}
\end{tcolorbox}

\begin{resultbox}[Conclusion de la phase]
La preuve est nette : le couple \texttt{(username, authenticationDatabase)} fait partie intégrante de l'identité MongoDB. La compréhension de ce point est démontrée opérationnellement.
\end{resultbox}

\subsection{Phase 5 -- Construction de la politique RBAC}

\subsubsection{Création du rôle personnalisé \texttt{incidentManager}}

L'étape suivante consiste à créer un rôle personnalisé limité à la collection \texttt{incidents} :

\begin{tcolorbox}[breakable,colback=CodeBg,colframe=CodeFrame,title=Création du rôle personnalisé]
\begin{Verbatim}[fontsize=\small,breaklines,breakanywhere]
use soc_lite
db.createRole({
  role: "incidentManager",
  privileges: [
    {
      resource: { db: "soc_lite", collection: "incidents" },
      actions: ["find", "insert", "update", "remove"]
    }
  ],
  roles: []
})
\end{Verbatim}
\end{tcolorbox}

Cette étape est structurante pour la suite : elle matérialise le passage d'un modèle de droits générique à un modèle de droits finement ciblé.

\subsubsection{Insertion des données de test}

Les captures montrent ensuite l'insertion d'un jeu de données cohérent pour pouvoir tester les droits :
\begin{itemize}
  \item trois événements dans \texttt{security\_logs} (échecs et succès SSH),
  \item deux incidents dans \texttt{incidents},
  \item deux utilisateurs dans \texttt{users}.
\end{itemize}

Les objets visibles incluent notamment :
\begin{itemize}
  \item l'incident \texttt{INC-2026-0042} : brute-force SSH sur \texttt{srv-web-01},
  \item l'incident \texttt{INC-2026-0041} : scan de ports depuis \texttt{192.168.5.200},
  \item les utilisateurs \texttt{analyst\_dupont} et \texttt{analyst\_martin}.
\end{itemize}

\subsubsection{Création des comptes applicatifs et métiers}

Les comptes suivants sont observés :
\begin{itemize}
  \item \texttt{soc\_analyst\_ro} avec le rôle built-in \texttt{read} sur \texttt{soc\_lite},
  \item \texttt{soc\_manager} avec le rôle personnalisé \texttt{incidentManager},
  \item \texttt{soc\_app} avec le rôle built-in \texttt{readWrite} sur \texttt{soc\_lite}.
\end{itemize}

Les commandes sont visibles dans les captures et confirmées par \texttt{db.getUsers()} et \texttt{db.getRoles()}.

\subsection{Phase 6 -- Tests d'accès par profil}

\subsubsection{\texttt{soc\_analyst\_ro} : lecture autorisée, écriture refusée, mais lecture de \texttt{users} possible}

Les preuves montrent la séquence suivante :
\begin{itemize}
  \item lecture des logs : \textbf{autorisée},
  \item lecture des incidents de sévérité \texttt{high} : \textbf{autorisée},
  \item insertion dans \texttt{security\_logs} : \textbf{refusée},
  \item suppression d'un incident : \textbf{refusée},
  \item lecture de \texttt{users} : \textbf{autorisée en pratique}.
\end{itemize}

\begin{warnbox}[Observation critique]
Le support pédagogique visait une politique où l'analyste n'accède pas à \texttt{users}. Or la capture montre que \texttt{db.users.find()} retourne bien les deux documents. Ce comportement est cohérent avec MongoDB : le rôle built-in \texttt{read} sur une base s'applique à \textbf{toutes les collections} de cette base. Il n'est donc pas suffisamment restrictif pour satisfaire le moindre privilège sur \texttt{soc\_lite}.
\end{warnbox}

\subsubsection{\texttt{soc\_manager} : rôle personnalisé correctement borné}

La séquence pour \texttt{soc\_manager} est la plus conforme à la politique cible :
\begin{itemize}
  \item lecture des incidents : \textbf{autorisée},
  \item mise à jour de \texttt{INC-2026-0042} : \textbf{autorisée} (\texttt{modifiedCount: 1}),
  \item lecture des logs : \textbf{refusée},
  \item lecture de \texttt{users} : \textbf{refusée},
  \item commande d'administration \texttt{listDatabases} : \textbf{refusée}.
\end{itemize}

\begin{resultbox}[Évaluation]
Le rôle personnalisé \texttt{incidentManager} réalise effectivement une séparation propre des privilèges : accès en lecture/écriture sur \texttt{incidents} uniquement, sans débordement vers les logs, les utilisateurs ou l'administration du serveur.
\end{resultbox}

\subsubsection{\texttt{soc\_app} : fonctionnement applicatif validé, mais privilèges trop larges}

La séquence observée pour \texttt{soc\_app} montre :
\begin{itemize}
  \item insertion d'un nouvel événement \texttt{firewall\_drop} dans \texttt{security\_logs} : \textbf{autorisée},
  \item lecture des incidents ouverts : \textbf{autorisée},
  \item lecture de \texttt{users} : \textbf{autorisée en pratique},
  \item commande \texttt{listDatabases} : \textbf{refusée}.
\end{itemize}
\newpage
\begin{warnbox}[Point de sécurité restant]
Le compte \texttt{soc\_app} ne peut pas administrer l'instance, ce qui est positif. En revanche, le rôle built-in \texttt{readWrite} sur \texttt{soc\_lite} lui donne aussi accès à \texttt{users}. Le fonctionnement applicatif est donc validé, mais la politique de moindre privilège n'est pas totalement atteinte.
\end{warnbox}

\section{Analyse ``avant / après'' du durcissement}

\subsection{Avant durcissement}

\begin{longtable}{p{4.4cm}p{9.8cm}}
\toprule
\textbf{Élément} & \textbf{État observé} \\
\midrule
Authentification & Absente ; connexion anonyme possible \\
Port MongoDB & Publié sans restriction via \texttt{-p 27017:27017} \\
Lecture de données & Libre, sans identifiants \\
Écriture de données & Libre, sans identifiants \\
Suppression / destruction & Possible via \texttt{drop()} et \texttt{dropDatabase()} \\
Administration & Possible anonymement via \texttt{listDatabases} \\
\bottomrule
\end{longtable}

\subsection{Après durcissement}

\begin{longtable}{p{4.4cm}p{9.8cm}}
\toprule
\textbf{Élément} & \textbf{État observé} \\
\midrule
Authentification & Activée ; accès anonyme refusé \\
Port MongoDB & Restreint à \texttt{127.0.0.1} côté publication Docker \\
Compte root & Fonctionnel pour l'administration \\
authDB & Testée avec succès et échec contrôlé \\
RBAC personnalisé & Mis en place pour \texttt{soc\_manager} \\
Séparation applicative & Partielle : \texttt{soc\_manager} bien borné, \texttt{soc\_analyst\_ro} et \texttt{soc\_app} trop larges via rôles built-in \\
Administration par comptes non-root & Refusée \\
\bottomrule
\end{longtable}

\section{Checklist de durcissement et état d'avancement}

\rowcolors{2}{GrayLight}{white}
\begin{longtable}{p{6cm}C{2.4cm}p{5.8cm}}
\toprule
\textbf{Contrôle} & \textbf{État} & \textbf{Justification fondée sur les preuves} \\
\midrule
Démonstration du risque MongoDB sans auth & \statusok & Lecture, insertion, suppression de collection et \texttt{dropDatabase()} observés sans authentification \\
Activation de l'authentification & \statusok & Refus explicite des commandes anonymes \\
Création d'un compte administrateur initial & \statusok & Compte root validé sur instance sécurisée et instance manuelle \\
Validation de \texttt{authenticationDatabase} & \statusok & Succès sur \texttt{soc\_lite} et échec sur \texttt{admin} pour \texttt{soc\_app} \\
Restriction de l'exposition réseau & \statusok & Publication Docker liée à \texttt{127.0.0.1} \\
Création d'un rôle personnalisé fin & \statusok & \texttt{incidentManager} créé et testé \\
Création d'un profil analyste en lecture & \statuspartial & Lecture et refus d'écriture validés, mais accès à \texttt{users} encore possible \\
Création d'un profil manager limité aux incidents & \statusok & Accès incidents validé, logs/users/admin refusés \\
Création d'un profil applicatif limité & \statuspartial & Fonction applicative validée, mais accès à \texttt{users} encore possible \\
Gestion des secrets via \texttt{.env} / externalisation & \statusno & Aucune preuve fournie dans les captures \\
Checklist signée & \statusno & Non visible dans les éléments transmis \\
Devoir maison de correction finale des rôles built-in & \statusno & Pas de preuve d'une remédiation effective pour exclure \texttt{users} \\
\bottomrule
\end{longtable}

\section{Évaluation finale}

\subsection{Ce qui est pleinement réussi}

Le travail remis montre une maîtrise solide des fondamentaux opérationnels de sécurité MongoDB :
\begin{itemize}
  \item l'étudiant a \textbf{prouvé} le danger d'une instance sans authentification ;
  \item il a \textbf{basculé vers une instance protégée} et vérifié le refus d'accès anonyme ;
  \item il a \textbf{créé et utilisé} un administrateur initial ;
  \item il a \textbf{validé la notion d'authDB} avec un test positif et un test négatif ;
  \item il a \textbf{construit un rôle personnalisé pertinent} pour la gestion des incidents ;
  \item il a \textbf{testé systématiquement} les opérations autorisées et refusées selon chaque profil.
\end{itemize}

\subsection{Ce qui reste à finaliser pour un moindre privilège strict}

Le point principal restant concerne la collection \texttt{users}. Les rôles built-in \texttt{read} et \texttt{readWrite} sont trop larges pour satisfaire la politique métier visée. Pour achever le devoir de durcissement, il faudrait :
\begin{enumerate}
  \item remplacer \texttt{read} pour \texttt{soc\_analyst\_ro} par un rôle personnalisé limitant la lecture à \texttt{security\_logs} et \texttt{incidents} ;
  \item remplacer \texttt{readWrite} pour \texttt{soc\_app} par un rôle personnalisé limitant l'écriture à \texttt{security\_logs} et, si nécessaire, la lecture ciblée de \texttt{incidents} ;
  \item conserver \texttt{users} hors du périmètre des comptes applicatifs et métiers ;
  \item externaliser les mots de passe dans un \texttt{.env} ou un gestionnaire de secrets au lieu de les laisser visibles en ligne de commande.
\end{enumerate}

\begin{warnbox}[Formulation honnête de conclusion]
Sur la base des preuves fournies, le \textbf{TP principal est réalisé avec succès}. La partie \textbf{devoir maison} est en revanche \textbf{seulement préparée} par l'observation du problème, mais \textbf{non démontrée comme effectivement terminée} dans les captures reçues.
\end{warnbox}

\section{Conclusion}

Ce rapport met en évidence une progression nette entre un état initial \emph{insecure} et un état partiellement durci, avec authentification, compte administrateur, contrôle d'accès RBAC et tests de validation. La réalisation est techniquement sérieuse et pédagogique, car elle ne se limite pas à ``faire marcher'' les commandes : elle vérifie les effets attendus et les refus de sécurité.

Le résultat le plus important, du point de vue cybersécurité, est aussi le plus instructif : les rôles built-in semblent pratiques, mais ils ne suffisent pas toujours à respecter le \emph{least privilege}. Le travail fourni le démontre concrètement. C'est précisément ce type de constat qui transforme un TP en apprentissage de sécurité utile.

\appendix
\clearpage
\section{Annexe A -- Journal de preuves visuelles}

Cette annexe reprend l'intégralité des captures fournies et les annote de manière synthétique. L'objectif est de conserver la traçabilité des manipulations réellement exécutées.

\evidencepage{fig01.png}{Capture 1 -- Instance \texttt{mongo\_insecure} et destruction anonyme de la base.}{La capture montre le lancement d'un conteneur MongoDB sans authentification, la connexion anonyme via \texttt{mongosh}, puis des opérations destructrices acceptées : \texttt{db.incidents.drop()} retourne \texttt{true} et \texttt{db.dropDatabase()} supprime \texttt{soc\_lite}. C'est la preuve la plus forte du risque initial.}

\evidencepage{fig02.png}{Capture 2 -- Insertion anonyme et commande d'administration sur l'instance insecure.}{Un document \texttt{test\_insecure} est inséré sans identifiants dans \texttt{security\_logs}, puis retrouvé avec succès. La commande \texttt{db.adminCommand(\{ listDatabases: 1 \})} réussit également. Cela établit la compromission totale possible en mode non authentifié.}

\evidencepage{fig03.png}{Capture 3 -- Passage à \texttt{mongo\_secure} avec publication locale et refus d'accès anonyme.}{Le conteneur insecure est supprimé puis remplacé par \texttt{mongo\_secure}, lié à \texttt{127.0.0.1}. Une tentative de lecture de \texttt{security\_logs} échoue avec \texttt{Command find requires authentication}.}

\evidencepage{fig04.png}{Capture 4 -- Validation du compte administrateur sur l'instance sécurisée.}{Après authentification en tant qu'administrateur, \texttt{listDatabases} fonctionne. Cela confirme que la sécurité distingue correctement un client non authentifié d'un client autorisé.}

\evidencepage{fig05.png}{Capture 5 -- Création manuelle du compte root sur \texttt{mongo\_manual}.}{Cette capture documente un scénario plus avancé : démarrage manuel avec \texttt{mongod --auth --bind\_ip\_all}, création du compte root dans la base \texttt{admin}, puis reconnexion réussie.}

\evidencepage{fig06.png}{Capture 6 -- Authentification réussie de \texttt{soc\_app} avec \texttt{authenticationDatabase = soc\_lite}.}{Le compte applicatif se connecte avec la bonne base d'authentification et peut interroger \texttt{security\_logs}. La capture sert de validation positive du mécanisme d'authentification par base.}

\evidencepage{fig07.png}{Capture 7 -- Échec d'authentification de \texttt{soc\_app} avec la mauvaise authDB.}{Ici, la même identité applicative est utilisée contre \texttt{admin}. MongoDB retourne \texttt{Authentication failed}. Cette capture est importante car elle prouve que le paramètre d'authDB n'est pas décoratif : il conditionne réellement l'accès.}

\evidencepage{fig08.png}{Capture 8 -- Création du rôle personnalisé \texttt{incidentManager}.}{La commande \texttt{db.createRole()} définit un rôle limité à la collection \texttt{incidents} avec les actions \texttt{find}, \texttt{insert}, \texttt{update} et \texttt{remove}. C'est la base du profil manager.}

\evidencepage{fig09.png}{Capture 9 -- Insertion du jeu de données de test dans \texttt{security\_logs}, \texttt{incidents} et \texttt{users}.}{La capture montre l'insertion de plusieurs journaux SSH, de deux incidents et de deux utilisateurs. Elle fournit le socle de données nécessaire aux vérifications RBAC ultérieures.}

\evidencepage{fig10.png}{Capture 10 -- Création de \texttt{soc\_analyst\_ro} et \texttt{soc\_manager}.}{Le compte analyste reçoit le rôle built-in \texttt{read} sur \texttt{soc\_lite}, tandis que le compte manager reçoit le rôle personnalisé \texttt{incidentManager}.}

\evidencepage{fig11.png}{Capture 12 -- Vérification des utilisateurs et rôles présents dans \texttt{soc\_lite}.}{La sortie de \texttt{db.getUsers()} confirme la présence de \texttt{soc\_analyst\_ro}, \texttt{soc\_app} et \texttt{soc\_manager}. La sortie de \texttt{db.getRoles()} confirme la présence du rôle personnalisé \texttt{incidentManager}.}

\evidencepage{fig12.png}{Capture 13 -- Visualisation des rôles built-in et du rôle personnalisé.}{La commande \texttt{db.getRoles(\{showBuiltinRoles:true\})} rappelle que \texttt{read} et \texttt{readWrite} sont des rôles globaux à l'échelle de la base. Cette capture aide à expliquer les débordements de privilèges observés ensuite.}

\evidencepage{fig13.png}{Capture 14 -- Tests du compte \texttt{soc\_analyst\_ro}.}{Les lectures de logs et d'incidents sont autorisées, l'insertion et la suppression sont refusées, mais \texttt{db.users.find()} retourne les documents. La capture est centrale car elle matérialise l'écart entre la politique souhaitée et le comportement réel du rôle built-in \texttt{read}.}

\evidencepage{fig14.png}{Capture 15 -- Tests du compte \texttt{soc\_manager}.}{Le compte manager lit les incidents, modifie le statut d'un incident avec succès, mais ne peut ni lire \texttt{security\_logs}, ni lire \texttt{users}, ni administrer l'instance. C'est la mise en œuvre la plus propre du moindre privilège dans l'ensemble des preuves.}

\evidencepage{fig15.png}{Capture 16 -- Tests du compte \texttt{soc\_app}.}{Le compte applicatif insère un événement \texttt{firewall\_drop} et lit les incidents ouverts. En revanche, il peut aussi lire \texttt{users}, ce qui révèle la limite du rôle built-in \texttt{readWrite}. L'administration de l'instance reste heureusement refusée.}


\end{document}
